\documentclass[11pt,a4paper]{article}
\usepackage[utf8]{inputenc}
\usepackage[T1]{fontenc}
\usepackage[norsk]{babel}
\usepackage{tgtermes}
\usepackage{tgheros}
\renewcommand{\ttdefault}{txtt}
\usepackage{amsmath,amssymb}
\usepackage[a4paper,left=2.45cm,right=2.6cm,top=2.3cm,bottom=2.2cm,headsep=0.9cm]{geometry}
\usepackage{fancyhdr}
\usepackage{lastpage}
\usepackage{xcolor}
\usepackage{tikz}
\usepackage{pgfplots}
\pgfplotsset{compat=1.18}
\usepackage{booktabs}
\usepackage{array}
\usepackage{enumitem}
\usepackage{listings}
\usepackage{needspace}
\usepackage{multirow}
\usetikzlibrary{positioning,arrows.meta,calc}
\definecolor{insgrey}{gray}{0.42}
\definecolor{radiogrey}{gray}{0.70}
\setlength{\parindent}{0pt}
\setlength{\parskip}{0.45em}
\setlist[itemize]{itemsep=0pt,topsep=2pt,label=\textbullet,font=\sffamily}
\lstset{basicstyle=\ttfamily\small,columns=fullflexible,keepspaces=true,
  showstringspaces=false,upquote=true,aboveskip=0.6em,belowskip=0.6em,
  literate={æ}{{\ae}}1 {ø}{{\o}}1 {å}{{\aa}}1 {Æ}{{\AE}}1 {Ø}{{\O}}1 {Å}{{\AA}}1}
\pagestyle{fancy}
\fancyhf{}
\renewcommand{\headrulewidth}{0pt}
\fancyhead[R]{\fontsize{7.5}{9}\selectfont\sffamily\color{insgrey}TDT4172 - Introduksjon til maskinlæring - Øvingssett 10}
\fancyfoot[R]{\fontsize{7}{9}\selectfont\sffamily\color{insgrey}\thepage/\pageref{LastPage}}
\newcommand{\delstart}[1]{\clearpage
  \noindent\hspace*{-1.55em}\makebox[1.55em][l]{\sffamily\bfseries\fontsize{12}{14}\selectfont\raisebox{0.35em}{#1}}%
  {\sffamily\bfseries\fontsize{18}{22}\selectfont Part #1}\par\vspace{1.0em}}
\newcommand{\sporsmal}[1]{\par\Needspace{9\baselineskip}\vspace{1.9em}%
  \noindent{\bfseries Spørsmål #1}\par\vspace{-0.45em}%
  \noindent{\bfseries -{}-{}-{}-{}-{}-{}-{}-{}-{}-{}-{}-{}-{}-{}-{}-{}-}\par\vspace{-0.15em}}
\newcommand{\velg}[1]{\par\Needspace{6\baselineskip}\noindent{\bfseries #1}\par\vspace{0.3em}}
\newcommand{\radio}[1]{\tikz[baseline=-0.62ex]{\draw[radiogrey,line width=0.45pt,fill=gray!5] (0,0) circle (0.56em);
  \node[inner sep=0pt,font=\sffamily\fontsize{5.5}{6}\selectfont,text=insgrey] at (0,0) {#1};}}
\newenvironment{alternativer}{\begin{list}{}{\setlength{\leftmargin}{2.35em}\setlength{\labelwidth}{1.4em}%
  \setlength{\labelsep}{0.6em}\setlength{\itemsep}{0.85em}\setlength{\topsep}{0.45em}\setlength{\parsep}{0pt}}%
  \sffamily\fontsize{10}{12.6}\selectfont}{\end{list}}
\newcommand{\alt}[2]{\item[\radio{#1}]#2}
\newcommand{\makspoeng}{\par\vspace{2.2em}\hfill{\sffamily\fontsize{8.5}{10}\selectfont Maks poeng: 10}\hspace*{0.5em}\par}
\begin{document}

\noindent\hspace*{-1.55em}\makebox[1.55em][l]{\raisebox{0.3em}{\sffamily\bfseries\fontsize{11}{12}\selectfont i}}%
{\sffamily\bfseries\fontsize{18}{22}\selectfont Frontpage}\par\vspace{1.1em}

{\bfseries Eksamensoppgave i TDT4172 - Introduksjon til maskinlæring}\\
{\bfseries Øvingssett:} 10 av 10 (uoffisielt øvingssett)\\
{\bfseries Tid:} 4 timer

\vspace{0.6em}
{\bfseries Faglig kontakt:} --- (øvingssett, ikke en offisiell eksamen)\\
{\bfseries Møter i eksamenslokalet:} Nei

\vspace{0.6em}
{\bfseries Hjelpemiddelkode/Tillatte hjelpemidler:} D / Ingen trykte eller håndskrevne hjelpemidler tillatt.
Bestemt, enkel kalkulator tillatt.

\vspace{0.6em}
{\bfseries\underline{ANNEN INFORMASJON:}}

\vspace{0.4em}
{\bfseries Les oppgavene nøye og gjør dine egne antagelser.} Presiser i besvarelsen hvilke forutsetninger du
har lagt til grunn i tolkning/avgrensing av oppgaven.

Dette øvingssettet er laget for egentrening, med samme format og vanskelighetsgrad som de
siste ordinære eksamenene i emnet. Løsningsforslaget står på siste side.

\vspace{1.6em}
{\bfseries\underline{FAGSPESIFIKK INFORMASJON}}

\vspace{0.4em}
{\bfseries Ingen papirhåndtegninger:}\\
Denne eksamenen tillater ikke bruk av håndtegninger.

\vspace{1.4em}
{\bfseries Vekting av oppgavene}:\\
Denne eksamenen inneholder flervalgsoppgaver organisert i ti grupper på fire spørsmål hver. Hver
gruppe vurderes separat etter følgende regler:

\vspace{0.6em}
For hver av de fire spørsmålene i gruppen får du følgende poeng
\begin{itemize}[leftmargin=2.2em,label=\textperiodcentered,itemsep=0pt]
\item {\bfseries Riktig svar:} +2,5 poeng
\item {\bfseries Feil svar (et alternativ er valgt, men ikke korrekt):} $-$0,83 poeng
\item {\bfseries Ikke besvart:} 0 poeng
\end{itemize}
{\bfseries Merk} at dersom du først har valgt et svaralternativ, kan du {\bfseries ikke velge det bort} igjen. Da må du
svare på det spørsmålet. Med andre ord mister du muligheten til å la et spørsmål stå blankt så fort
du har trykket på et alternativ.

\vspace{0.6em}
Poengsummen for hver gruppe er {\bfseries summen} av poengene for de fire enkeltspørsmålene.\\
For {\bfseries alle spørsmålene} gjelder en {\bfseries null-grense}: dersom summen for denne gruppen blir negativ,
settes verdien til {\bfseries 0}. Den totale poengsummen på eksamen er summen av poengene fra alle
gruppene.

\delstart{1}
\sporsmal{1}
En logistisk regresjonsmodell gir $z = -6$ for et treningspunkt med $y = 1$, altså
$\hat y = \sigma(-6) \approx 0.0025$: modellen er svært sikker, men tar feil. Med cross-entropy loss er
$\partial L/\partial w_j = (\hat y - y)\,x_j$. Tenk deg at du i stedet bruker $L = \tfrac12(\hat y - y)^2$ med
den samme sigmoid-modellen, og husk at $\sigma'(z) = \sigma(z)\,(1-\sigma(z))$.

Hvordan blir gradienten for $w_j$ fra dette punktet med de to tapsfunksjonene?
\velg{Velg ett alternativ:}
\begin{alternativer}
\alt{A}{Med begge tapene er den nær $-x_j$, siden valget av tapsfunksjon ikke påvirker hvor raskt feil rettes}
\alt{B}{Med begge tapene er den nesten null, siden sigmoid er mettet og punktet derfor bidrar lite uansett}
\alt{C}{Med cross-entropy er den nesten null fordi sigmoid er mettet, mens kvadratisk tap gir en gradient nær $-x_j$}
\alt{D}{Med cross-entropy er den nær $-x_j$, mens faktoren $\hat y(1-\hat y)$ gjør den nesten null med kvadratisk tap}
\end{alternativer}
\sporsmal{2}
Et sykehus trener en modell som skal predikere om en innlagt pasient får diagnosen sepsis; omtrent
10\,\% av pasientene i dataene får det. En av featurene er \emph{Antall\_doser\_sepsismedisin}, hentet
fra journalen ved utskrivning. Modellen får 99\,\% accuracy og recall $0.97$ på et tilfeldig utvalgt
valideringssett, men gjør det svært dårlig når den tas i bruk ved innleggelse.

Hva er den mest sannsynlige årsaken?
\velg{Velg ett alternativ:}
\begin{alternativer}
\alt{A}{Modellen er overtilpasset treningsdataene, siden så høy accuracy er urealistisk for en diagnose}
\alt{B}{Featuren registreres etter at sepsis er påvist, så den lekker target og er ukjent ved bruk}
\alt{C}{Modellen predikerer nesten bare den største klassen, slik at accuracy blir høy men verdiløs}
\alt{D}{Featuren er sterkt korrelert med andre features, så vektene blir ustabile når nye data kommer}
\end{alternativer}
\sporsmal{3}
En logistisk regresjonsmodell med parametre $\mathbf{w}, b$ gir $\hat y = (0.9,\ 0.2,\ 0.6)$ for tre
treningspunkter med $y = (1,\ 0,\ 1)$. Når dataene er gitt og uttrykket
$\prod_i \hat y_i^{\,y_i}(1-\hat y_i)^{1-y_i}$ betraktes som funksjon av $\mathbf{w}$ og $b$, kalles
det \textbf{likelihood}.

Hva er \textbf{log-likelihood} (med $\ln$) for disse dataene, og hva gjør treningen med den?
\velg{Velg ett alternativ:}
\begin{alternativer}
\alt{A}{Den er $-0.84$, og treningen justerer $\mathbf{w}, b$ slik at den blir så stor som mulig}
\alt{B}{Den er $-0.84$, og treningen justerer $\mathbf{w}, b$ slik at den blir så liten som mulig}
\alt{C}{Den er $-2.23$, og treningen justerer $\mathbf{w}, b$ slik at den blir så stor som mulig}
\alt{D}{Den er $-2.23$, og treningen justerer $\mathbf{w}, b$ slik at den blir så liten som mulig}
\end{alternativer}
\sporsmal{4}
Du trener en logistisk regresjonsmodell med mini-batch stochastic gradient descent.

Hvilken liste inneholder bare \textbf{hyperparametre}?
\velg{Velg ett alternativ:}
\begin{alternativer}
\alt{A}{Læringsraten $\eta$, antall epoker og batch size}
\alt{B}{Biasen $b$, læringsraten $\eta$ og batch size}
\alt{C}{Vektene $\mathbf{w}$, biasen $b$ og batch size}
\alt{D}{Læringsraten $\eta$, antall epoker og vektene $\mathbf{w}$}
\end{alternativer}
\makspoeng

\delstart{2}
\sporsmal{1}
En bildeklassifikator for tre klasser gir følgende confusion matrix på testsettet. Radene er
\textbf{sann} klasse og kolonnene er \textbf{predikert} klasse.
\begin{center}
\begin{tabular}{lccc}
\toprule
 & Pred.\ katt & Pred.\ hund & Pred.\ fugl \\
\midrule
Sann katt & 40 & 8  & 2  \\
Sann hund & 10 & 30 & 0  \\
Sann fugl & 5  & 7  & 18 \\
\bottomrule
\end{tabular}
\end{center}
Hva er \textbf{precision} for klassen \emph{hund}?
\velg{Velg ett alternativ:}
\begin{alternativer}
\alt{A}{$0.79$}
\alt{B}{$0.75$}
\alt{C}{$0.67$}
\alt{D}{$0.73$}
\end{alternativer}
\sporsmal{2}
Valideringssettet har 50 positive punkter. Tabellen viser antall true positives og false positives
for tre klassifiseringsterskler.
\begin{center}
\begin{tabular}{lcc}
\toprule
Terskel $t$ & TP & FP \\
\midrule
0.3 & 48 & 60 \\
0.5 & 42 & 18 \\
0.7 & 33 & 6 \\
\bottomrule
\end{tabular}
\end{center}
Kravet er recall $\ge 0.80$, og blant tersklene som oppfyller kravet vil du ha høyest mulig precision.
Hvilken terskel velges, og hva er precision der?
\velg{Velg ett alternativ:}
\begin{alternativer}
\alt{A}{$t = 0.7$, med precision $0.85$}
\alt{B}{$t = 0.5$, med precision $0.70$}
\alt{C}{$t = 0.5$, med precision $0.84$}
\alt{D}{$t = 0.3$, med precision $0.44$}
\end{alternativer}
\sporsmal{3}
I et datasett for svindeldeteksjon er klasse 1 (svindel) sterkt underrepresentert. En kollega bruker
den vektede tapsfunksjonen
\[ L = -w_1\, y \ln \hat y - w_0\,(1-y)\ln(1-\hat y), \]
men setter ved en feil $w_0 = 10$ og $w_1 = 1$.

Hva er den mest sannsynlige effekten, sammenlignet med å trene med uvektet tap?
\velg{Velg ett alternativ:}
\begin{alternativer}
\alt{A}{Feil på klasse 0 straffes hardest, så modellen predikerer svindel enda sjeldnere og recall faller typisk}
\alt{B}{Vektingen virker som å senke terskelen til 0.1, så modellen gir flere alarmer og precision faller}
\alt{C}{Vektene endrer bare tapsverdien og ikke gradientene, så den ferdig trente modellen blir den samme}
\alt{D}{Feil på klasse 1 straffes hardest, så modellen predikerer svindel oftere og recall øker typisk}
\end{alternativer}
\sporsmal{4}
Et testsett utvides med 5\,000 ekstra negative punkter. Alle er svært lette, slik at modellen
klassifiserer samtlige som negative, og ingen nye falske positive oppstår. Modellen og terskelen er uendret.

Hvordan påvirker dette metrikkene som rapporteres for den positive klassen?
\velg{Velg ett alternativ:}
\begin{alternativer}
\alt{A}{Accuracy, precision og F1 øker, mens recall og FPR er uendret}
\alt{B}{Alle metrikkene endres, siden andelen positive i testsettet blir lavere}
\alt{C}{Precision og accuracy øker, mens recall, FPR og F1 er uendret}
\alt{D}{Accuracy og TNR øker og FPR synker, mens precision, recall og F1 er uendret}
\end{alternativer}
\makspoeng

\delstart{3}
\sporsmal{1}
Tabellen viser antall passasjerer i Titanic-data fordelt på klasse og overlevelse.
\begin{center}
\begin{tabular}{lcc}
\toprule
 & \emph{Survived}=0 & \emph{Survived}=1 \\
\midrule
\emph{Pclass}=1 & 80  & 136 \\
\emph{Pclass}=2 & 97  & 87  \\
\emph{Pclass}=3 & 372 & 119 \\
\bottomrule
\end{tabular}
\end{center}
Hva er $p(\text{\emph{Pclass}}=1 \mid \text{\emph{Survived}}=1)$?
\velg{Velg ett alternativ:}
\begin{alternativer}
\alt{A}{$0.40$}
\alt{B}{$0.24$}
\alt{C}{$0.38$}
\alt{D}{$0.15$}
\end{alternativer}
\sporsmal{2}
Volumet av en $d$-dimensjonal kule med radius $r$ er $K_d\, r^d$, der $K_d$ bare avhenger av $d$.
Se på en 50-dimensjonal kule med radius 1.

Hvor stor andel av volumet ligger i det ytre skallet med tykkelse $\epsilon = 0.05$, altså mellom radius 0.95 og 1?
\velg{Velg ett alternativ:}
\begin{alternativer}
\alt{A}{$0.05$}
\alt{B}{$0.08$}
\alt{C}{$0.14$}
\alt{D}{$0.92$}
\end{alternativer}
\sporsmal{3}
En kollega bytter \texttt{criterion='gini'} til \texttt{criterion='entropy'} i
\texttt{DecisionTreeClassifier} og trener på nytt med de samme dataene og
\texttt{max\_depth=6}.

Hva er det mest sannsynlige resultatet?
\velg{Velg ett alternativ:}
\begin{alternativer}
\alt{A}{Treningen går som regel raskere, fordi entropi er standardvalget i sklearn og derfor optimalisert}
\alt{B}{Treet blir som regel mye dypere, men like raskt å trene, fordi entropi har maksimum 1 mens Gini har 0.5}
\alt{C}{Treet blir som regel helt annerledes, fordi Gini og entropi rangerer splittene i motsatt rekkefølge}
\alt{D}{Treet blir som regel nesten likt, men treningen går litt tregere fordi entropi krever logaritmer}
\end{alternativer}
\sporsmal{4}
Et beslutningstre bygges rekursivt etter pseudokoden fra forelesningen, med
\texttt{max\_depth=5}. En node på dybde 2 inneholder 3 positive og 2 negative datapunkter, og alle
features er allerede brukt på veien ned fra rotnoden. Foreldrenoden hadde 4 positive og 9 negative.

Hva skjer med noden?
\velg{Velg ett alternativ:}
\begin{alternativer}
\alt{A}{Den blir en løvnode som predikerer klasse 1, siden det ikke finnes flere features å splitte på}
\alt{B}{Den blir en løvnode som predikerer klasse 0, siden den dominante klassen i foreldrenoden brukes}
\alt{C}{Den fjernes, og datapunktene sendes tilbake til foreldrenoden for å prøve en annen splitt}
\alt{D}{Den splittes videre på en tilfeldig valgt feature, siden maksimal dybde ikke er nådd ennå}
\end{alternativer}
\makspoeng

\delstart{4}
\sporsmal{1}
Modell A predikerer alltid gjennomsnittet av target i treningsdataene. Modell B er et regresjonstre
uten dybdebegrensning. Begge trenes på mange ulike treningssett fra samme fordeling, og du ser på
prediksjonene i ett fast punkt $x$ der den sanne sammenhengen er langt fra gjennomsnittet.

Hvilket utsagn er riktig?
\velg{Velg ett alternativ:}
\begin{alternativer}
\alt{A}{A har lav varians og høy bias, mens B har lav bias og høy varians}
\alt{B}{A har høy varians og høy bias, mens B har lav bias og lav varians}
\alt{C}{A har høy varians og lav bias, mens B har høy bias og lav varians}
\alt{D}{A har lav varians og lav bias, mens B har høy bias og høy varians}
\end{alternativer}
\sporsmal{2}
Du har 10 treningspunkter med ulike $x$-verdier og målestøy i $y$. Du tilpasser et polynom av grad 9
med vanlig minste kvadraters metode, uten regularisering.

Hva er mest sannsynlig?
\velg{Velg ett alternativ:}
\begin{alternativer}
\alt{A}{Trenings-MSE blir 0 bare hvis dataene er uten støy, og ellers omtrent like stor som valideringsfeilen}
\alt{B}{Trenings-MSE blir 0, siden 10 parametre alltid kan gå eksakt gjennom 10 punkter, mens valideringsfeilen blir høy}
\alt{C}{Trenings-MSE blir høy, siden et polynom av grad 9 er for stivt for 10 punkter, og valideringsfeilen blir lav}
\alt{D}{Trenings-MSE blir omtrent lik støyvariansen, siden støy ikke kan tilpasses bort, og valideringsfeilen blir lav}
\end{alternativer}
\sporsmal{3}
I et Titanic-prosjekt gjør du tre ting: (i) du dropper kolonnen \emph{PassengerId}, (ii) du skalerer
\emph{Age} til intervallet $[0,1]$, og (iii) du lager en ny feature $\text{\emph{Age}}^2$.

Hvilken type feature engineering er henholdsvis (i), (ii) og (iii)?
\Needspace{9\baselineskip}
\velg{Velg ett alternativ:}
\begin{alternativer}
\alt{A}{(i) feature preprocessing, (ii) feature extraction, (iii) feature selection}
\alt{B}{(i) feature selection, (ii) feature preprocessing, (iii) feature extraction}
\alt{C}{(i) feature preprocessing, (ii) feature selection, (iii) feature extraction}
\alt{D}{(i) feature extraction, (ii) feature selection, (iii) feature preprocessing}
\end{alternativer}
\sporsmal{4}
En regresjonsmodell har MSE $= 900$ (uten faktoren $\tfrac12$) på valideringsdata. Targetene ligger
typisk rundt $28$.

Hvilken tolkning er riktig?
\velg{Velg ett alternativ:}
\begin{alternativer}
\alt{A}{MSE under 1000 regnes generelt som lavt, så modellen er god uansett hva targetene er}
\alt{B}{Typisk avvik er rundt 3, omtrent en tidel av targetene, så modellen er ganske god}
\alt{C}{Typisk avvik er rundt 30, omtrent like stort som targetene selv, så modellen er dårlig}
\alt{D}{Typisk avvik er rundt 900, langt større enn targetene, så modellen er ubrukelig}
\end{alternativer}
\makspoeng

\delstart{5}
\sporsmal{1}
I en random forest har hvert tre prediksjonsvarians $\sigma^2 = 4$, og parvis korrelasjon mellom
trærnes prediksjoner er $\rho = 0.6$. Variansen til gjennomsnittet av $B$ trær er
\[ \rho\,\sigma^2 + \frac{1-\rho}{B}\,\sigma^2 . \]

Hvilken verdi nærmer variansen til ensemblet seg når antall trær $B$ blir svært stort?
\velg{Velg ett alternativ:}
\begin{alternativer}
\alt{A}{$0.00$}
\alt{B}{$2.40$}
\alt{C}{$1.60$}
\alt{D}{$0.08$}
\end{alternativer}
\sporsmal{2}
Etter første runde i et boosting-ensemble for regresjon har punkt $i$ target $10$ og prediksjon $7$,
mens punkt $j$ har target $4$ og prediksjon $4.5$.

Hvordan skiller \textbf{AdaBoost} og \textbf{gradient boosting} seg i neste runde?
\velg{Velg ett alternativ:}
\begin{alternativer}
\alt{A}{Begge tilpasser neste tre til $3$ og $-0.5$, men AdaBoost trekker i tillegg et nytt bootstrap-utvalg}
\alt{B}{AdaBoost tilpasser neste tre til $3$ og $-0.5$, mens gradient boosting øker vekten til punktene med størst feil}
\alt{C}{Begge øker vekten til punkt $i$, men gradient boosting erstatter i tillegg targets med residualene}
\alt{D}{AdaBoost øker vekten til punktene med størst feil, mens gradient boosting tilpasser neste tre til $3$ og $-0.5$}
\end{alternativer}
\sporsmal{3}
Du skal trene 400 trær, og hvert tre tar 2 sekunder å trene på én CPU-kjerne. Maskinen har 8 kjerner.
Anta at parallellisering bare kan skje ved å trene ulike trær samtidig.

Hvor lang tid tar treningen omtrent for en \textbf{random forest} og for \textbf{gradient boosting}?
\velg{Velg ett alternativ:}
\begin{alternativer}
\alt{A}{Random forest 800 s og gradient boosting 800 s}
\alt{B}{Random forest 100 s og gradient boosting 100 s}
\alt{C}{Random forest 800 s og gradient boosting 100 s}
\alt{D}{Random forest 100 s og gradient boosting 800 s}
\end{alternativer}
\sporsmal{4}
En bank har brukt ett beslutningstre med dybde 3 til å vurdere lånesøknader, og saksbehandlerne kan følge
reglene fra rotnoden til en løvnode. Treet byttes ut med en \textbf{random forest} med 500 fullt utvokste
trær, som gir høyere accuracy på testdata.

Hva er den viktigste ulempen med byttet?
\velg{Velg ett alternativ:}
\begin{alternativer}
\alt{A}{Skogen gir bare en klasse og ikke en sannsynlighet, så saksbehandlerne får ikke vite hvor sikker den er}
\alt{B}{Hver prediksjon er en avstemning over 500 ulike regelsett, så den kan ikke leses som én kjede av splitter}
\alt{C}{Fullt utvokste trær gjør at skogens prediksjoner varierer mer enn det ene treets når treningsdata endres}
\alt{D}{Skogen kan ikke brukes på nye søkere før alle de 500 trærne er trent på nytt med de nye søkernes data}
\end{alternativer}
\makspoeng

\delstart{6}
\sporsmal{1}
Output-laget i et nettverk har tre noder med \textbf{softmax}, og tapet er \textbf{cross-entropy}
$L = -\sum_k y_k \ln \hat y_k$. For denne kombinasjonen er $\partial L / \partial z_k = \hat y_k - y_k$, der
$z_k$ er logiten (verdien før softmax) til node $k$. For et datapunkt med target $\mathbf{y} = (0,\ 1,\ 0)$
gir nettverket $\hat{\mathbf{y}} = (0.1,\ 0.2,\ 0.7)$.

Hvordan flytter ett gradient descent-steg logitene for dette datapunktet?
\velg{Velg ett alternativ:}
\begin{alternativer}
\alt{A}{$z_3$ opp, $z_2$ kraftig ned og $z_1$ litt ned}
\alt{B}{$z_2$ opp, mens $z_1$ og $z_3$ ikke endres}
\alt{C}{$z_2$ opp, $z_3$ kraftig ned og $z_1$ litt ned}
\alt{D}{$z_2$ opp, $z_1$ kraftig ned og $z_3$ litt ned}
\end{alternativer}
\sporsmal{2}
Et output-lag med softmax får logits $\mathbf{z} = (2,\ 1,\ 0)$ for ett datapunkt og
$\mathbf{z}' = (12,\ 11,\ 10)$ for et annet.

Hvilket utsagn om de predikerte sannsynlighetene er riktig?
\velg{Velg ett alternativ:}
\begin{alternativer}
\alt{A}{De kan ikke sammenlignes, siden logitene til det andre punktet ikke ligger i intervallet $[0, 1]$}
\alt{B}{Det andre punktet får jevnere sannsynligheter, siden forskjellene er små relativt til logitene}
\alt{C}{De er like for begge punktene, siden softmax bare avhenger av forskjellene mellom logitene}
\alt{D}{Det andre punktet får nesten all sannsynlighet på første klasse, siden logitene er mye større}
\end{alternativer}
\sporsmal{3}
Et nettverk har én skjult ReLU-node og én lineær output:
\[ h = \mathrm{ReLU}(w_1 x + b_1), \qquad \hat y = w_2 h + b_2, \qquad L = \tfrac12(\hat y - y)^2 , \]
med $w_1 = -1$, $b_1 = 0.5$, $w_2 = 3$, $b_2 = 1$. For datapunktet $x = 2$, $y = 0$ beregnes gradientene
med backpropagation.

Hvilke parametre får gradient null for dette datapunktet?
\velg{Velg ett alternativ:}
\begin{alternativer}
\alt{A}{$w_1$, $b_1$, $w_2$ og $b_2$}
\alt{B}{$w_1$, $b_1$ og $w_2$}
\alt{C}{Bare $w_1$ og $b_1$}
\alt{D}{Bare $w_2$}
\end{alternativer}
\sporsmal{4}
Et perceptron har vektene $\mathbf{w} = (2,\ 1,\ 1)$ og output $H(\mathbf{w}\mathbf{x} + b)$, der $H(z) = 1$ hvis
$z \ge 0$ og ellers 0. Det skal være aktivt nøyaktig når $2x_1 + x_2 + x_3 \ge 3$.

Hvilken verdi må biasen $b$ ha?
\velg{Velg ett alternativ:}
\begin{alternativer}
\alt{A}{$b = -1.5$}
\alt{B}{$b = 3$}
\alt{C}{$b = -3$}
\alt{D}{$b = -4$}
\end{alternativer}
\makspoeng

\delstart{7}
\sporsmal{1}
Du kjører k-means med $k = 4$ fem ganger på samme data, med ulike tilfeldige startsentroider.
Inertia etter konvergens blir $412,\ 388,\ 388,\ 530$ og $395$.

Hva bør du gjøre?
\velg{Velg ett alternativ:}
\begin{alternativer}
\alt{A}{Øke til $k = 5$ og kjøre på nytt, siden ulik inertia mellom kjøringene viser at $k = 4$ er feil}
\alt{B}{Velge en kjøring med inertia 388, siden alle har samme $k$ og lav inertia betyr kompakte klynger}
\alt{C}{Velge kjøringen med inertia 530, siden de andre kjøringene trolig har overtilpasset seg dataene}
\alt{D}{Lete etter en feil i koden, siden k-means skal konvergere til samme løsning fra hvilken som helst start}
\end{alternativer}
\sporsmal{2}
Du har kjørt k-means for $k = 2, \dots, 6$ og beregnet tre metrikker:
\begin{center}
\begin{tabular}{lccccc}
\toprule
$k$ & 2 & 3 & 4 & 5 & 6 \\
\midrule
Inertia & 980 & 610 & 450 & 390 & 340 \\
Silhouette & 0.48 & 0.61 & 0.55 & 0.42 & 0.39 \\
Calinski--Harabasz & 310 & 420 & 465 & 400 & 360 \\
\bottomrule
\end{tabular}
\end{center}
Hvilken $k$ gir \textbf{best verdi} for henholdsvis silhouette, Calinski--Harabasz og inertia?
\velg{Velg ett alternativ:}
\begin{alternativer}
\alt{A}{Silhouette $k=6$, Calinski--Harabasz $k=4$ og inertia $k=6$}
\alt{B}{Silhouette $k=3$, Calinski--Harabasz $k=4$ og inertia $k=6$}
\alt{C}{Silhouette $k=6$, Calinski--Harabasz $k=2$ og inertia $k=2$}
\alt{D}{Silhouette $k=3$, Calinski--Harabasz $k=2$ og inertia $k=6$}
\end{alternativer}
\sporsmal{3}
Kunder beskrives med \emph{Inntekt} (kr) og \emph{Alder} (år). Kunden $\mathbf{p} = (500\,000,\ 30)$ skal
tilordnes en av sentroidene $\mathbf{c}_A = (505\,000,\ 60)$ og $\mathbf{c}_B = (480\,000,\ 31)$ med
euklidsk avstand. Standardavviket er $100\,000$ kr for inntekt og $10$ år for alder.

Hvilken sentroide tilordnes kunden uten skalering, og hvilken etter standardisering av begge
featurene?
\velg{Velg ett alternativ:}
\begin{alternativer}
\alt{A}{Uten skalering $\mathbf{c}_A$, etter standardisering $\mathbf{c}_A$}
\alt{B}{Uten skalering $\mathbf{c}_B$, etter standardisering $\mathbf{c}_B$}
\alt{C}{Uten skalering $\mathbf{c}_A$, etter standardisering $\mathbf{c}_B$}
\alt{D}{Uten skalering $\mathbf{c}_B$, etter standardisering $\mathbf{c}_A$}
\end{alternativer}
\sporsmal{4}
En klynge i et endimensjonalt datasett består av punktene $2,\ 3,\ 4,\ 5,\ 6$. I tillegg blir en
uteligger i $x = 40$ tilordnet den samme klyngen, fordi den er nærmere denne sentroiden enn noen annen.

Hvor havner sentroiden etter neste oppdatering i k-means?
\velg{Velg ett alternativ:}
\begin{alternativer}
\alt{A}{$10.0$}
\alt{B}{$21.0$}
\alt{C}{$4.5$}
\alt{D}{$4.0$}
\end{alternativer}
\makspoeng

\delstart{8}
\sporsmal{1}
Du har bare 6 datapunkter, hver med 50 features, og kjører PCA (dataene sentreres først).

Hvor mange hovedkomponenter kan høyst ha egenverdi større enn 0?
\velg{Velg ett alternativ:}
\begin{alternativer}
\alt{A}{$5$}
\alt{B}{$44$}
\alt{C}{$50$}
\alt{D}{$6$}
\end{alternativer}
\sporsmal{2}
Et isolasjonstre splitter en node ved å velge en tilfeldig feature og en splittverdi $s$ trukket
uniformt mellom minste og største verdi av featuren i noden. Punkter med $x \le s$ går til venstre,
resten til høyre. En node inneholder ti punkter med én feature:
\[ 1,\ 2,\ 3,\ 4,\ 5,\ 6,\ 7,\ 8,\ 9,\ 30 . \]
Hva er sannsynligheten for at punktet 30 havner alene i en løvnode allerede etter første splitt?
\velg{Velg ett alternativ:}
\begin{alternativer}
\alt{A}{$0.28$}
\alt{B}{$0.72$}
\alt{C}{$0.70$}
\alt{D}{$0.10$}
\end{alternativer}
\sporsmal{3}
I t-SNE er perplexity lik $2^{H}$, der $H = -\sum_j p_{j|i}\log_2 p_{j|i}$ er entropien til
naboskapsfordelingen rundt punktet $\mathbf{x}_i$. Anta at $p_{j|i} = 0.25$ for fire punkter og
tilnærmet 0 for alle andre.

Hva blir $H$ og perplexity?
\velg{Velg ett alternativ:}
\begin{alternativer}
\alt{A}{$H = 2.00$ og perplexity $4.00$}
\alt{B}{$H = 1.39$ og perplexity $2.61$}
\alt{C}{$H = -2.00$ og perplexity $0.25$}
\alt{D}{$H = 0.50$ og perplexity $1.41$}
\end{alternativer}
\sporsmal{4}
En kollega bygger bare ett isolasjonstre i stedet for en Isolation Forest med 100 trær. Hvilke punkter
som får kortest sti til rotnoden, endrer seg mye fra kjøring til kjøring.

Hva er den beste forklaringen på at en skog med mange trær fungerer bedre?
\velg{Velg ett alternativ:}
\begin{alternativer}
\alt{A}{Hvert nytt tre legger mer vekt på punktene de forrige trærne isolerte tidlig, så feilene rettes gradvis opp}
\alt{B}{Hvert tre stemmer «anomali» eller «normal», og et punkt regnes som anomali når flertallet av trærne er enige}
\alt{C}{Hvert tre splitter tilfeldig, så stilengden i ett enkelt tre er støyete, mens snittet over mange trær er stabilt}
\alt{D}{Snittet over mange trær fjerner en systematisk skjevhet i stilengden, så anomaliene får kortere sti enn i ett tre}
\end{alternativer}
\makspoeng

\delstart{9}
\sporsmal{1}
En robotstøvsuger får reward $+1$ for hver enhet støv den suger opp. Etter trening har den lært å
tømme beholderen på gulvet og suge opp det samme støvet igjen, om og om igjen.

Hva er den beste forklaringen og løsningen?
\velg{Velg ett alternativ:}
\begin{alternativer}
\alt{A}{Q-verdiene har ikke konvergert ennå, så agenten bør trenes i flere episoder til den finner en bedre policy}
\alt{B}{Agenten har utforsket for lite til å oppdage at et rent rom lønner seg; $\epsilon$ bør derfor økes igjen}
\alt{C}{$\gamma$ er for liten, så agenten ser ikke verdien av å rydde hele rommet på sikt; $\gamma$ bør derfor økes mot 1}
\alt{D}{Rewarden måler bare en tilnærming til målet, og agenten maksimerer den slik den står; den bør omformuleres}
\end{alternativer}
\sporsmal{2}
En episode gir rewardene $R_1 = 0,\ R_2 = 0,\ R_3 = 1,\ R_4 = 0,\ R_5 = 5$, der $R_{t+1}$ er rewarden
agenten får etter handlingen i steg $t$. Episoden slutter etter $R_5$, og diskonteringsfaktoren er
$\gamma = 0.9$.

Hva er den diskonterte returen $G_0 = \sum_{k=0}^{4} \gamma^k R_{k+1}$?
\velg{Velg ett alternativ:}
\begin{alternativer}
\alt{A}{$5.81$}
\alt{B}{$5.05$}
\alt{C}{$6.00$}
\alt{D}{$4.09$}
\end{alternativer}
\sporsmal{3}
En agent bruker $\epsilon$-greedy med $\epsilon = 0.2$ i en tilstand med 5 mulige handlinger. I
tilfeldig modus trekkes handlingen uniformt blant \emph{alle} 5 handlinger, også den grådige.

Hva er sannsynligheten for at agenten velger den grådige handlingen (høyest Q-verdi)?
\velg{Velg ett alternativ:}
\begin{alternativer}
\alt{A}{$0.96$}
\alt{B}{$0.20$}
\alt{C}{$0.84$}
\alt{D}{$0.80$}
\end{alternativer}
\sporsmal{4}
En agent med Q-tabell trenes i et rutenett, alltid fra den samme startruten, til den følger korteste vei
til målet. Når den grådige agenten etterpå settes i andre ruter, når den ofte ikke målet. Trenes den
i stedet med en tilfeldig startrute i hver episode, når den målet fra alle ruter.

Hva er den beste forklaringen?
\velg{Velg ett alternativ:}
\begin{alternativer}
\alt{A}{$\gamma$ er for lav, så verdien av målet når ikke frem til ruter som ligger langt fra startruten}
\alt{B}{Q-tabellen generaliserer fra besøkte ruter til nabo-rutene, men bare innenfor et begrenset område}
\alt{C}{Q-tabellen har overtilpasset seg den ene veien, og bør regulariseres slik et nevralt nettverk gjøres}
\alt{D}{Q-verdiene oppdateres bare for tilstander som besøkes, og mange ruter ble sjelden eller aldri besøkt}
\end{alternativer}
\makspoeng

\delstart{10}
\sporsmal{1}
En modell er gjennomsnittet av to trær, $f(\mathbf{x}) = \tfrac12\bigl(f_1(\mathbf{x}) + f_2(\mathbf{x})\bigr)$, og alle
SHAP-verdier beregnes med samme bakgrunnsdatasett. For et datapunkt har featuren \emph{Alder} SHAP-verdi
$0.20$ i tre 1 og $0.08$ i tre 2.

Hva er SHAP-verdien til Alder for hele modellen, og hvilket aksiom gir dette?
\velg{Velg ett alternativ:}
\begin{alternativer}
\alt{A}{$0.20$, siden symmetri gjør at feature-verdien bestemmes av treet der den bidrar mest}
\alt{B}{$0.12$, siden effektivitet gjør at bidragene fra de to trærne må trekkes fra hverandre}
\alt{C}{$0.28$, siden additivitet gjør at verdiene for et snitt av modeller er summen av verdiene}
\alt{D}{$0.14$, siden additivitet gjør at verdiene for et snitt av modeller er snittet av verdiene}
\end{alternativer}
\sporsmal{2}
Et team skal forklare et nevralt nettverk for lånesøknader til saksbehandlere. Forslag 1 er en
fullstendig utskrift av alle vekter og aktiveringer. Forslag 2 er et beslutningstre med to regler som
gir samme avgjørelse (innvilg/avslå) som nettverket i 55\,\% av søknadene.

Hvilken vurdering er riktig etter kursets krav til en forklaring?
\velg{Velg ett alternativ:}
\begin{alternativer}
\alt{A}{Begge er gode, siden begge er post-hoc og kan lages uten å endre det ferdig trente nettverket}
\alt{B}{Utskriften er best, siden det viktigste kravet til en forklaring er at den gjengir modellen helt nøyaktig}
\alt{C}{Ingen av dem: utskriften er treffsikker men ikke forståelig, og treet er forståelig men ikke treffsikkert}
\alt{D}{Treet er best, siden det viktigste kravet til en forklaring er at mennesker kan forstå og bruke den}
\end{alternativer}
\sporsmal{3}
Se på koden under.
\par\noindent\begin{minipage}{\linewidth}
\begin{lstlisting}[language=Python]
probs = model.predict(X_test)   # softmax-output med shape (500, 3)
y_pred = ...                    # linjen som mangler
\end{lstlisting}
\end{minipage}
Hvilken linje gir én predikert klasse (0, 1 eller 2) for hvert av de 500 testpunktene?
\velg{Velg ett alternativ:}
\begin{alternativer}
\alt{A}{\texttt{y\_pred = np.argmax(probs, axis=0)}}
\alt{B}{\texttt{y\_pred = np.max(probs, axis=1)}}
\alt{C}{\texttt{y\_pred = np.argmax(probs, axis=1)}}
\alt{D}{\texttt{y\_pred = (probs > 0.5).astype(int)}}
\end{alternativer}
\sporsmal{4}
Tilstandene i et miljø med 16 ruter one-hot-enkodes før de sendes inn i et Q-nettverk:
\par\noindent\begin{minipage}{\linewidth}
\begin{lstlisting}[language=Python]
import numpy as np
states = np.array([3, 0, 15])
onehot = np.eye(16)[states]
print(onehot.shape, onehot.sum(axis=1), np.argmax(onehot, axis=1))
\end{lstlisting}
\end{minipage}
Hvilken beskrivelse av utskriften er riktig?
\velg{Velg ett alternativ:}
\begin{alternativer}
\alt{A}{Shape $(3, 16)$, radsummene er $1, 1, 1$, og argmax gir $3, 0, 15$}
\alt{B}{Shape $(3,)$, radsummene er $16, 16, 16$, og argmax gir $3, 0, 15$}
\alt{C}{Shape $(16, 3)$, radsummene er $1, 1, 1$, og argmax gir $3, 0, 15$}
\alt{D}{Shape $(3, 16)$, radsummene er $3, 0, 15$, og argmax gir $1, 1, 1$}
\end{alternativer}
\makspoeng

\clearpage
{\sffamily\bfseries\fontsize{18}{22}\selectfont Løsningsforslag}\par\vspace{0.3em}
\noindent\rule{\textwidth}{0.4pt}\par\vspace{1.2em}
Riktig alternativ per spørsmål (bokstaven står i rundingen foran hvert alternativ).\par\vspace{1em}
\begin{center}\sffamily\renewcommand{\arraystretch}{1.35}
\begin{tabular}{l*{4}{>{\centering\arraybackslash}p{1.6cm}}}
\toprule
& \textbf{Sp.\ 1} & \textbf{Sp.\ 2} & \textbf{Sp.\ 3} & \textbf{Sp.\ 4}\\
\midrule
Part 1 & D & B & A & A\\
Part 2 & C & B & A & D\\
Part 3 & A & D & D & A\\
Part 4 & A & B & B & C\\
Part 5 & B & D & D & B\\
Part 6 & C & C & B & C\\
Part 7 & B & B & C & A\\
Part 8 & A & B & A & C\\
Part 9 & D & D & C & D\\
Part 10 & D & C & C & A\\
\bottomrule
\end{tabular}
\end{center}
\end{document}
